\section*{第 \ref{ch:b1:topology} 章 --- 实直线上的拓扑}

\begin{solution}{exo:b1:topology:1}
$\intoo{0}{1} \cup \intoo{2}{3}$：开集（开集之并），不
是闭集（$\frac 1n \to 0$ 落在集合之外）。

$\intco{0}{1}$：既不是开集（集合内没有包含 $0$ 的区间），
也不是闭集（$1 - \frac1n \to 1 \notin$ 集合）。

$\{0\} \cup \intcc{1}{2}$：闭集（有限个闭集之并），不
是开集（在 $0$ 处不满足要求）。

$\R \setminus \Z$：开集（$\Z$ 为闭集），不是闭集：数列
$\bigl(\frac 1n\bigr)$ 落在该集合中，但其极限 $0$ 属于 $\Z$，
即其极限逃出了该集合。

$\Q \cap \intoo{0}{1}$：既非开集（有理数的任一邻域都含无理数），
也非闭集：它含有趋于无理数 $\frac{\sqrt 2}{2}$ 的数列，
由稠密性可知。
\end{solution}

\begin{solution}{exo:b1:topology:2}
$A = \intoc{0}{1} \cup \{2\}$：$\mathring A = \intoo{0}{1}$，
$\overline A = \intcc{0}{1} \cup \{2\}$，$\partial A = \{0, 1, 2\}$。

$A = \R \setminus \Q$：$\mathring A = \emptyset$（每个区间
都含有有理数），$\overline A = \R$（无理数稠密），
$\partial A = \R$。

$A = \bigl\{\frac{(-1)^n n}{n+1}\bigr\}$：偶数项趋于 $1$，
奇数项趋于 $-1$，二者均不属于 $A$。$\mathring A =
\emptyset$（孤立点），$\overline A = A \cup \{-1, 1\}$，
$\partial A = \overline A$。
\end{solution}

\begin{solution}{exo:b1:topology:3}
\emph{补集。} 对 $F = \{a_1 < a_2 < \dots < a_k\}$，其补集
是开区间的并：$\intoo{-\infty}{a_1}$、
$\intoo{a_i}{a_{i+1}}$、$\intoo{a_k}{+\infty}$ --- 由
\cref{prop:b1:topology:openstable}可知其为开集。

\emph{数列。} 设 $u_n \in F$，$u_n \to \ell$。取 $\varepsilon
= \min\{\abs{a_i - a_j} : i \neq j\}/2 > 0$（或任取 $\varepsilon$，当
$F$ 为单点集时）：自某项起，各项与极限的距离
都在 $\varepsilon$ 之内（即到 $\ell$ 的距离），从而两两之间都在 $2\varepsilon$ 之内，
这迫使它们从该项起都等于同一个 $a_i$；于是
$\ell = a_i \in F$。
\end{solution}

\begin{solution}{exo:b1:topology:4}
$\subseteq$：$\overline A \cup \overline B$ 是闭集（有限并），
且包含 $A \cup B$，故它包含\emph{最小的}闭
超集 $\overline{A \cup B}$。$\supseteq$：$A \subseteq A \cup B$
给出 $\overline A \subseteq \overline{A \cup B}$（闭包运算是
单调的：$A$ 的附着点也是较大集合的附着点），而
对 $B$ 同理。

交集的反例：$A = \intoo{0}{1}$，$B =
\intoo{1}{2}$：$\overline{A \cap B} = \overline\emptyset = \emptyset$
而 $\overline A \cap \overline B = \{1\}$。
\end{solution}

\begin{solution}{exo:b1:topology:5}
利用闭集的数列刻画
（见\cref{thm:b1:topology:seqclosed}）。设 $S = \{u_n\} \cup \{\ell\}$
，$(v_k)$ 为 $S$ 中的数列，$v_k \to m$；需证 $m \in S$。分两种
情形。若某个值 $v \in S$ 被 $(v_k)$ 无穷多次取到，
则常数子列给出 $m = v \in S$。否则每个
值都只被取有限次；特别地，对每个 $n$，项 $u_n$
只出现有限次，$\ell$ 亦然。于是对每个 $N$，
指标 $k$（满足 $v_k \in \{u_0, \dots, u_N, \ell\}$）只有有限多个，
其余的 $v_k$ 均为项 $u_n$，其中 $n > N$。给定
$\varepsilon > 0$，取 $N$ 使得 $\abs{u_n - \ell} \leq
\varepsilon$ 对所有 $n > N$ 成立，则除有限多个 $v_k$ 外都满足
$\abs{v_k - \ell} \leq \varepsilon$。于是 $v_k \to \ell$，故 $m =
\ell \in S$。
\end{solution}

\begin{solution}{exo:b1:topology:6}
$U + A = \bigcup_{a \in A} (U + a)$，且每个平移 $U + a$
均为开集（平移各区间即可验证）。开集之并
仍为开集（见\cref{prop:b1:topology:openstable}）。

闭集：$\Z$ 与 $\sqrt 2\,\Z$ 均为闭集（$\alpha\Z$ 中的
收敛数列最终为常数，参见
\cref{ex:b1:topology:closed}）。二者之和 $\Z + \sqrt 2\,\Z$ 在
$\R$ 中稠密（见\cref{exo:b1:reals:9}），但不等于 $\R$（它是可数的，或
更简单地：$\frac{\sqrt 2}{2} \notin \Z + \sqrt2\Z$，否则 $\sqrt 2$
将是有理数 --- 设 $\frac{\sqrt2}{2} = m + n\sqrt 2$
推出 $(2n - 1)\sqrt 2 = -2m$，故 $\sqrt 2 \in \Q$，除非 $n =
\frac12$（不可能）。稠密的真子集不是闭集：其
闭包为 $\R \neq$ 自身。
\end{solution}

\begin{solution}{exo:b1:topology:7}
$\Z$：$\intoo{n - \frac12}{n + \frac12}$ 为 $n$ 的邻域
仅与 $\Z$ 相交于 $n$。

$A = \{\frac1n : n \in \N^*\}$：在 $\frac 1n$ 附近，半径为
$\frac{1}{n} - \frac{1}{n+1} = \frac{1}{n(n+1)}$ 的区间（例如取一半）
可将其与邻近点隔开 --- 因而该集合的各点都是孤立点。但 $0 \in
\overline A \setminus A$：孤立点并不排除附着的
外点。

若 $A$ 的所有点都是孤立点，则 $A$ 中没有内点，因为
内点周围有一个完整的区间落在 $A$ 中
（区间为无穷集），与孤立性矛盾。故 $\mathring A
= \emptyset$。
\end{solution}

\begin{solution}{exo:b1:topology:8}
若 $A \subseteq \intcc{-M}{M}$，闭集 $\intcc{-M}{M}$
包含 $A$，从而包含 $\overline A$（最小的闭
超集）：$\overline A$ 有界。

设 $s = \sup A$（有限）。由于 $A \subseteq \overline A$，$\sup
\overline A \geq s$。反之，$\overline A \subseteq
\intoc{-\infty}{s}$：该半直线为闭集，且包含 $A$，故
$\overline A$ 的每个元素都 $\leq s$，从而 $\sup\overline A
\leq s$。于是等号成立。

$s \in \overline A$：由 $\varepsilon$ 刻画
（\cref{prop:b1:reals:epsilon}），每个区间 $\intoo{s -
\varepsilon}{s + \varepsilon}$ 都含 $A$ 中的元素，故 $s$ 为
附着点。
\end{solution}

\begin{solution}{exo:b1:topology:9}
设 $A$ 既开又闭，$a \in A$，$b \in \R
\setminus A$，不妨设 $a < b$。集合 $B = A \cap
\intcc{a}{b}$ 非空（含 $a$）且有界，令 $s = \sup B$。由
\cref{exo:b1:topology:8}，$s \in \overline B \subseteq \overline A =
A$（$A$ 为闭集）。注意 $s \leq b$，又 $b \notin A$，故 $s < b$。
由 $A$ 为开集，存在区间 $\intoo{s - r}{s + r}$ 含于 $A$，
且可取 $r < b - s$。于是 $s + \frac{r}{2}$ 属于 $A
\cap \intcc{a}{b} = B$ 且大于 $s$ --- 与 $s = \sup
B$ 矛盾。因此不存在这样的 $(a, b)$：$A$、$\R \setminus A$ 中
之一为空集。
\end{solution}

\begin{solution}{exo:b1:topology:10}
$I_x$ 为若干个都含 $x$ 的开区间之并：它是开集，
又因其凸，故为区间 --- 若 $u < z < v$，$u, v \in
I_x$，则 $u$、$v$ 分别位于开子区间 $J_u \ni x$、$J_v \ni
x$，它们都含于 $U$，且 $J_u \cup J_v$ 为区间（二者都含 $x$），
它含于 $U$ 且含 $z$，故 $z \in I_x$
（\cref{prop:b1:reals:intervals}）。

若 $I_x \cap I_y \neq \emptyset$，则 $I_x \cup I_y$ 为开区间
（凸性：两个相交区间之并），含于 $U$，
且含 $x$ 与 $y$，由极大性，$I_x \cup I_y \subseteq I_x$，且
$\subseteq I_y$（由各自的极大性），故 $I_x = I_y$。

因此 $U$ 是互不相交的不同集合 $I_x$ 之并（每个 $x \in
U$ 都属于其自身的 $I_x$）。可数性：族中每个非空开区间
$I$ 含有有理数 $q_I$
（见\cref{thm:b1:reals:density}），且不同的不交区间
对应不同的有理数：该族单射地映入可数集 $\Q$，
（可由整数对编号）。故
至多可数个区间。
\end{solution}

\begin{solution}{exo:b1:topology:11}
$\overline A = A \cup A'$。（$\supseteq$）$A \subseteq \overline
A$ 恒成立；若 $x \in A'$，则 $x$ 的每个邻域都与 $A
\setminus \{x\} \subseteq A$ 相交，故 $x$ 为附着点。（$\subseteq$）
设 $x \in \overline A$。若 $x \in A$，结论成立；若 $x \notin A$，
则 $x$ 的每个邻域都与 $A = A \setminus \{x\}$ 相交，故 $x \in
A'$。

由此 $A$ 为闭集 $\iff$ $A = \overline A = A \cup A'$
$\iff$ $A' \subseteq A$。

$A = \{\frac 1n\}$：$0$ 为聚点（$\frac 1n \to
0$，各项 $\neq 0$）；每个 $\frac 1n$ 都是孤立点
（见\cref{exo:b1:topology:7}），故不在 $A'$ 中；若 $x
\notin A \cup \{0\}$，则其有一个与 $A$ 不相交的完整区间（位于
$x$ 在 $A \cup \{0\}$ 中的两个邻点之间，或在 $1$ 之外）。
故 $A' = \{0\}$。

$\Z' = \emptyset$：每个整数都是孤立点，每个非整数
都有一个含于 $\R \setminus \Z$ 的邻域。

$\Q' = \R$：任一实数周围的每个区间都含无穷
多个有理数（见\cref{thm:b1:reals:density}），特别地，其中必有一个
异于中心的有理数。
\end{solution}

\begin{solution}{exo:b1:topology:12}
\begin{enumerate}
  \item 对每个 $a \in A$：$\abs{x - a} \leq \abs{x - y} +
        \abs{y - a}$，故 $d_A(x) \leq \abs{x - y} + \abs{y - a}$；
        对 $a$ 取下确界，得 $d_A(x) \leq \abs{x - y} +
        d_A(y)$。交换 $x$、$y$ 的位置即得另一个
        不等式：$\abs{d_A(x) - d_A(y)} \leq \abs{x - y}$。
  \item $d_A(x) = 0$ $\iff$ 对每个 $\varepsilon > 0$，存在
        $a \in A$ 使 $\abs{x - a} < \varepsilon$ $\iff$ 任何以
        $x$ 为中心的区间都与 $A$ 相交 $\iff$ $x \in \overline A$。
        若 $F$ 为闭集且 $x \notin F = \overline F$，则
        $d_F(x) \neq 0$，即 $d_F(x) > 0$。
  \item 若 $d_F(x) < \varepsilon$，令 $r = \varepsilon -
        d_F(x) > 0$。对 $\abs{y - x} < r$，由 (1) 得
        $d_F(y) \leq d_F(x) + \abs{x - y} < \varepsilon$，故
        区间 $\intoo{x - r}{x + r}$ 含于 $V_\varepsilon$，
        因此它是开集，它包含 $F$，因为其上 $d_F = 0$
        成立。最后，$x \in \bigcap_{\varepsilon>0}
        V_\varepsilon$ $\iff$ $d_F(x) < \varepsilon$ 对所有
        $\varepsilon$ $\iff$ $d_F(x) = 0$ $\iff$ $x \in \overline
        F = F$。由于 $\bigcap_{\varepsilon > 0} V_\varepsilon =
        \bigcap_{n \geq 1} V_{1/n}$，故每个闭集都是
        可数个开集的交。
\end{enumerate}
\end{solution}

\begin{solution}{pb:b1:topology:1}
\textbf{1.} $C_1 = \intcc{0}{\frac13} \cup \intcc{\frac23}{1}$
和
\[
C_2 = \intcc{0}{\tfrac19} \cup \intcc{\tfrac29}{\tfrac13}
\cup \intcc{\tfrac23}{\tfrac79} \cup \intcc{\tfrac89}{1} .
\]
用归纳法：若 $C_n$ 是 $2^n$ 个互不相交的闭线段之并，
每段长为 $3^{-n}$，删去每段的开中间三分之一，
则每个父段留下两个长为 $3^{-n-1}$ 的闭线段：
共 $2^{n+1}$ 个线段，两两不交（不同父段的子段
因父段本身不交而互相分离；同一父段的子段
之间由被删去的间隙隔开）。

\textbf{2.} 每个 $C_n$ 是有限个线段的并，因而是闭集；
$C = \bigcap C_n$ 为闭集之交，
故亦为闭集（见\cref{def:b1:topology:closed}）；有界（$\subseteq
\intcc{0}{1}$），由\cref{thm:b1:topology:compact}知紧致。
非空：$0$ 落在每个 $C_n$ 最左端的线段中。设
$a$ 为某线段 $S$ 的端点，该线段属于 $C_n$。对 $m \leq n$，
有 $a \in C_n \subseteq C_m$。对之后的各阶段：中间三分之一
的去除从不删去端点，且 $a$ 仍是
$S$ 的两个子线段之一（即与 $a$ 相接的那一个）的端点；由
归纳法知 $a \in C_m$ 对所有 $m \geq n$ 成立，故 $a \in C$。

\textbf{3.} $C_n$ 的总长度为 $2^n \cdot 3^{-n} =
(\frac23)^n \to 0$。给定 $\varepsilon > 0$，取 $n$ 使
$(\frac23)^n \leq \varepsilon$，则 $C \subseteq C_n$，它是
有限个线段的并，总长度 $\leq \varepsilon$。

\textbf{4.} 设 $I \subseteq C$ 是含两个不同
点的区间。对每个 $n$ 有 $I \subseteq C_n$，而 $I$ 为凸集，
必落在 $C_n$ 的\emph{单个}线段之中（若与两个线段相交，
则 $I$ 将含有其间空隙中的一点，
而该点在 $C_n$ 之外）。因此 $I$ 的长度 $\leq
3^{-n}$ 对一切 $n$ 成立，矛盾。故 $C$ 中
的区间只能是空集或单点；特别地，不存在
$\intoo{x-r}{x+r}$ 含于 $C$，即 $\mathring C = \emptyset$。
由于 $C$ 为闭集，$\overline C = C$ 的内部为空：$C$ 是
无处稠密的。

\textbf{5.} 记 $\varphi_0(x) = \frac x3$，$\varphi_2(x) =
\frac{2 + x}{3}$，它们是 $\intcc{0}{1}$ 上的递增仿射双射，
分别映到 $\intcc{0}{\frac13}$ 与 $\intcc{\frac23}{1}$ 上。断言：
$C_{n+1} = \varphi_0(C_n) \cup \varphi_2(C_n)$。当 $n = 0$ 时，
即第 1 题。用归纳法：递增仿射映射将线段的
中间三分之一映到像线段的中间三分之一，
因此去除中间三分之一的操作与 $\varphi_0$ 和
$\varphi_2$ 可交换；对 $C_{n+1} =
\varphi_0(C_n) \cup \varphi_2(C_n)$ 施行去除步骤，即得 $C_{n+2} =
\varphi_0(C_{n+1}) \cup \varphi_2(C_{n+1})$。对
$n$ 取交：当 $x \leq \frac13$ 时，$x \in C \iff x \in \varphi_0(C_n)$
对一切 $n$ 成立 $\iff 3x \in \bigcap C_n = C$；在
$\intcc{\frac23}{1}$ 上同理；且 $\intoo{\frac13}{\frac23}$
中的点都不在 $C_1$ 中。因此 $C = \varphi_0(C) \cup \varphi_2(C)$，
且二者互不相交。

\textbf{6.} 对 $n$ 用归纳法。$n = 0$ 的情形说明每个 $x \in
\intcc{0}{1}$ 都有编码，这即\cref{pb:b1:reals:1}
（第 9 题，对 $x < 1$；$1 = (0.\overline 2)_3$）。设
在秩 $n$ 时的等价性成立。若 $x \in C_{n+1}$，由第 5 题，$x =
\varphi_i(z)$，其中 $z \in C_n$，$i \in \{0, 2\}$；若 $(e_k)$
是 $z$ 的编码且前 $n$ 位不含 $1$，则 $(i, e_1,
e_2, \dots)$ 的部分和为 $\frac i3 + \frac13\sum_{k\leq m}
e_k 3^{-k} \to \varphi_i(z) = x$，故为 $x$ 的编码，且前 $n +
1$ 位不含 $1$。反之，若 $x$ 有编码 $(d_k)$，其中
$d_1, \dots, d_{n+1} \in \{0, 2\}$，则平移后的串 $(d_2,
d_3, \dots)$ 有某个值 $z \in \intcc{0}{1}$，其前 $n$
位不含 $1$，将部分和的计算反过来看，
得 $x = \varphi_{d_1}(z)$；由归纳假设 $z \in
C_n$，故由第 5 题 $x \in C_{n+1}$。最后：一个
完全不含 $1$ 的编码使 $x$ 属于每个 $C_n$，从而属于 $C$；反之，
若 $x \in C$，则对每个 $n$，$x$ 的至多两个编码之一
（\cref{pb:b1:reals:1}，第 11 题）的前 $n$ 位不含 $1$；
必有一个固定编码对任意大的 $n$ 成立（由两个编码间的抽屉原理），
而若某编码前 $n$ 位
不含 $1$，且对任意大的 $n$ 成立，则
它完全不含 $1$。

\textbf{7.} $0 = (0.\overline 0)_3$，$1 = (0.\overline 2)_3$，
$\frac13 = (0.0\overline{2})_3$（即 $(0.1)_3$ 的非正常孪生），
$\frac23 = (0.2\overline{0})_3$。几何级数求和：
\[
(0.\overline{02})_3 = \sum_{j\geq1} \frac{2}{9^{\,j}}
= \frac{2/9}{1 - 1/9} = \frac14 ,
\qquad
(0.\overline{20})_3 = \sum_{j\geq1} \frac{2}{3\cdot 9^{\,j-1}}
= \frac{2/3}{1 - 1/9} = \frac34 ,
\]
二者均不含 $1$：$\frac14, \frac34 \in C$。（无穷和
是部分和上确界的缩写，如\cref{pb:b1:reals:1}所述。）
$C_n$ 的线段端点形如 $m/3^n$
（归纳：子线段的端点或为父线段的端点，或与之相差
$3^{-n-1}$ 的整数倍）。若 $\frac14 =
\frac{m}{3^n}$，则 $3^n = 4m$，而 $4 \nmid 3^n$，不可能。
因此 $\frac14 \in C$，但它从不是任何线段的端点。

\textbf{8.} 设 $x$ 有两个不同的不含 $1$ 的编码。
有两个编码本身意味着（\cref{pb:b1:reals:1}，第 11 题，以 $3$ 为底）
$x = m/3^N \in \intoo{0}{1}$，且两个编码分别为：
一个是有限展开式，其最后一个非零位 $d_N \in \{1, 2\}$，
其后接 $0$，另一个是其孪生，$d_N - 1$ 位于第 $N$ 位，
其后接 $2$。若 $d_N = 1$，则前者含有 $1$；若 $d_N
= 2$，则孪生中 $d_N - 1 = 1$。无论如何，这一对中至多一个
不含 $1$，矛盾。故每个 $x \in C$ 恰有
一个不含 $1$ 的编码（存在性由第 6 题给出），且
不同的 $\{0,2\}$ 串取值不同。每个
$\{0,2\}$ 串的值在 $\intcc{0}{1}$ 中（部分和 $\leq
1$），且各前缀均不含 $1$，故其值属于每个 $C_n$，
即属于 $C$：取值映射是从
$\{0,2\}$ 串集到 $C$ 上的双射。

\textbf{9.} 设 $(d^{(k)})$ 为不含 $1$ 的、$x_k$ 的编码，
令 $e_k = 2 - d^{(k)}_k \in \{0, 2\}$：这是一个 $\{0,2\}$ 串，其
取值 $y$ 位于 $C$ 中，且以 $(e_k)$ 为其唯一的不含 $1$ 的编码
（第 8 题）。对每个 $k$，$y$ 与 $x_k$ 的编码
在第 $k$ 位不同，故 $y \neq x_k$：不存在 $\N^* \to C$ 的满射。
一个长度为零的不可数集：基数意义上“大”，
测度意义上“小” --- 二者同时成立。

\textbf{10.} 紧致性：无穷交的闭性
加有界性（第 2 题） --- 这是唯一一个不能由单一
包含关系得到的性质。长度为零：$C \subseteq
C_n$，其总长度为 $(\frac23)^n$（第 3 题）。无处
稠密：$C \subseteq C_n$ 迫使 $C$ 内的区间长度
$\leq 3^{-n}$（第 4 题）。不可数性完全不依赖
任何包含关系：它需要完整的交结构，
并编码在第 8 题的双射之中。

\textbf{11.} 将 $d_n$ 换为 $2 - d_n$：新串仍是
$\{0,2\}$-串，故其值 $x_n$ 属于 $C$；第 $n$ 位以后的部分
和相差恰为 $2\cdot3^{-n}$，因此
$\abs{x_n - x} = 2\cdot3^{-n}$。于是 $x_n \neq x$ 且 $x_n \to
x$：$C$ 中每一点都是 $C$ 中其他点的极限 --- $C$
是\emph{完备集}，没有孤立点。

\textbf{12.} 由第 5 题的归纳，$C_n$ 的各线段恰为
$\intcc{t}{t + 3^{-n}}$，其中 $t$ 取遍
长度为 $n$ 的 $\{0,2\}$-串的值。设 $x \in C$
编码为 $(d_k)$，其截断 $t_n$（数字依次为 $d_1 \dots d_n$
后补若干 $0$），因此是左端点，且 $0 \leq x - t_n
\leq 3^{-n}$：端点在 $C$ 中稠密。它们构成
$\{m/3^n : m, n\}$ 的子集。该集合以整数对为指标，故
可数（与\cref{exo:b1:topology:10}中的 $\Q$ 相同）。由于 $C$
不可数（第 9 题），$C$ 中除可数多个点外的所有点
都不是端点 --- $\frac14$（第 7 题）是这座冰山的
可见一角。

\textbf{13.} 取 $n$ 使 $3^{-n} < y - x$。$x, y \in
C_n$，且二者不能位于同一线段（长度 $3^{-n} <
y - x$）。两线段之间被挖去的间隙给出 $z$
满足 $x < z < y$ 且 $z \notin C_n \supseteq C$。因此 $C$ 中任意两点
都被补集隔开：$C$ 的凸子集只能是
单点集 --- $C$ 完全不连通。

\textbf{14.} 若 $(d_k)$ 是无 $1$ 编码，其值为 $x$，则串
$(2 - d_k)$ 仍是 $\{0,2\}$-串，其部分和
\[
\sum_{k=1}^{n} (2 - d_k)3^{-k} = (1 - 3^{-n}) - \sum_{k=1}^n
d_k 3^{-k} \longrightarrow 1 - x :
\]
故 $1 - x \in C$。于是 $1 - C \subseteq C$，将该映射
作用两次得 $1 - C = C$：康托尔集
关于 $\frac12$ 对称。

\textbf{15.} 部分和 $t_n \to x$ 与 $t'_n \to x'$
（单调递增数列收敛于其上确界，即
其值），故 $t_n + t'_n \to x + x'$，由
\cref{thm:b1:seq:operations}。设 $y \in \intcc{0}{1}$，其
编码为 $(e_k)$。取 $(a_k, b_k) = (0,0), (0,2), (2,2)$
依 $e_k = 0, 1, 2$ 而定：则 $a_k + b_k = 2e_k$，
串 $(a_k)$、$(b_k)$ 为 $\{0,2\}$-串，其值分别为 $x,
x' \in C$，且
\[
x + x' = \lim_n\,(t_n + t'_n) = \lim_n 2\sum_{k=1}^n e_k 3^{-k}
= 2y :
\]
于是每个 $y \in \intcc{0}{1}$ 都是 $C$ 中两点的中点。

\textbf{16.} 由第 15 题得 $\intcc{0}{2} = 2\,\intcc{0}{1}
\subseteq C + C$，又 $C + C \subseteq \intcc{0}{1} +
\intcc{0}{1} = \intcc{0}{2}$，故等号成立。再利用 $1 - C = C$，
\[
C - C = C + (C - 1) = (C + C) - 1 = \intcc{-1}{1} .
\]
具体例子：$1 = \frac14 + \frac34$，是
$C$ 中两个非端点元素之和。

\textbf{17.} 长度度量的是集合本身占据直线的
部分；它对 \emph{和} 构成的集合一无所知，该集合
是二参数族 $C \times C$ 在 $(x,
x') \mapsto x + x'$ 下的像 --- 两个数字串是
相互独立选取的，正是这种自由度填满了
$\intcc{0}{2}$。没有任何定理能用加项的长度
界定和集的长度，而 $C$ 恰好证明这样的界不存在。

\textbf{18.} 由\cref{pb:b1:reals:1}（第 18 题），$x$ 为有理数
当且仅当其真展开最终周期。
无 $1$ 编码（对 $x \in C$ 而言）或为该真展开，或为
某有限展开的异形孪生；有限串与其
孪生（最终恒为 $2$）均最终
周期，故无 $1$ 编码的周期性等价于
$x$ 的有理性。对 $\frac1{13}$ 作 $3$ 进制长除法
（$r_0 = 1$）：$3 = 13\cdot0 + 3$，$9 = 13\cdot0 + 9$，$27 =
13\cdot2 + 1$，余数回到 $1$，数字为
$\overline{002}$，故 $\frac1{13} = (0.\overline{002})_3$，
无 $1$ 编码且周期，即 $C$ 中的有理点。（验算：
$\frac{2/27}{1 - 1/27} = \frac{2}{26} = \frac1{13}$。）

\textbf{19.} 令 $d_k = 2$ 于三角数位置，即
$k = \frac{j(j+1)}{2}$，其余为 $0$ 的串是 $\{0,2\}$-串，
故其值 $x^*$ 属于 $C$（第
8 题）。其中 $2$ 有无穷多个，相邻两者之间的间隔
为 $j + 1 \to \infty$，故它不是最终周期的（若周期为 $T$，则
最终必有 $2$ 以间隔 $\leq T$ 出现，此即
\cref{pb:b1:reals:1}中的间隔增大论证，见第 20 题）；由
第 18 题，$x^* \notin \Q$。又由第 9 题及
$\Q$ 的可数性，$C$ 中除可数多个元素外均为无理数：
$x^*$ 是常态，而非例外。

\textbf{20.} 设 $y \in \intcc{0}{1}$，它有二进制编码
$(c_k)$，其中 $c_k \in \{0, 1\}$（见\cref{pb:b1:reals:1}，第
9 题，$2$ 进制；$y = 1$ 取全为 $1$ 的串）。于是 $(2c_k)$
为 $\{0,2\}$-串，其值 $x$ 属于 $C$，且 $h(x)$ 为
$(c_k)$ 的值，即 $y$：$h$ 将 $C$ 满射到
$\intcc{0}{1}$。若 $C$ 是某个从 $\N^*$ 出发的映射之像，
则与 $h$ 复合即可列出 $\intco{0}{1}$ 的全部元素，
这与\cref{pb:b1:reals:1}中的对角线定理相矛盾
（第 22 题）：$C$ 再次不可数。长度为零的集合
满射到整个线段。

\textbf{21.} 由第 5 题，$C$ 为
$\varphi_0(C)$ 与 $\varphi_2(C)$ 的不交并，二者各为
$\frac13 C$ 的缩放副本的平移。由可加性、伸缩性与平移
不变性可得
\[
L(C) = L(\varphi_0(C)) + L(\varphi_2(C))
= \tfrac13 L(C) + \tfrac13 L(C) = \tfrac23\,L(C),
\]
于是 $\frac13 L(C) = 0$，故 $L(C) = 0$。仅凭自相似性
便判 $C$ 长度为零 --- 第 3 题只是执行了
这一判决。

\textbf{22.} 长为 $l_n$ 的线段去掉一个中央区间，
其长为 $4^{-(n+1)}$，留下两条长为 $l_{n+1}
= \frac{l_n - 4^{-(n+1)}}{2}$；由 $l_0 = 1$ 出发，
归纳可得 $l_n = \frac{2^n + 1}{2\cdot4^n}$：事实上
$\frac12\Bigl(\frac{2^n+1}{2\cdot4^n} - \frac{1}{4^{n+1}}\Bigr)
= \frac{2(2^n + 1) - 1}{2\cdot4^{n+1}} = \frac{2^{n+1} +
1}{2\cdot4^{n+1}}$，且 $l_n > 0$ 恒成立：构造过程
永不中断。$K = \bigcap K_n$ 是闭集且有界，从而为紧致集；
$K$ 内的区间落在 $K_n$ 的某条线段内，其长为
$l_n \to 0$，故内部为空。去掉的总长为 $\sum_{n\geq0} 2^n
\cdot 4^{-(n+1)} = \frac14\sum_{n\geq0}\bigl(\frac12\bigr)^n =
\frac12$，而每个 $K_n$ 的总长为 $2^n l_n = \frac{2^n +
1}{2^{n+1}} > \frac12$。现设有限个开区间
之并为 $U \supseteq K$。由
\cref{ex:b1:topology:nested}的配套结论，$U \supseteq K_n$ 对某个 $n$ 成立；
在本阶段承认有限个区间之并上长度的可加性，
则覆盖区间的总长不小于
$K_n$ 的总长，它超过 $\frac12$。故 $K$ 无处稠密，但
不存在廉价的覆盖：拓扑意义上的“小”（无处稠密）与
度量意义上的“小”（长度为零）是真正不同的概念，
而 $K$ 将二者区分开来。

\textbf{23.} 可达性：设 $d = d(x, F)$，取 $a_k \in F$
使 $\abs{x - a_k} \leq d + \frac1k$：$a_k$ 有界，因此
由波尔扎诺--魏尔斯特拉斯定理（\cref{thm:b1:seq:bw}）可取出
$a_{\varphi(k)} \to a$，且 $a \in F$（$F$ 为闭集，
\cref{thm:b1:topology:seqclosed}），且 $\abs{x - a} = \lim
\abs{x - a_{\varphi(k)}} = d$。下面求最大值：若 $y \in C$，
则 $d(y, C) = 0$；否则 $y$ 落在某一阶段
$n \geq 1$ 去掉的缺口中，它是长为 $3^{-n}$ 的开区间，其两端点
都属于 $C$（第 2 题），因此 $d(y, C) \leq
\frac{3^{-n}}{2} \leq \frac16$，等号成立要求 $n = 1$
且 $y$ 位于缺口 $\intoo{\frac13}{\frac23}$ 的中点，
即 $y = \frac12$；事实上 $d\bigl(\frac12, C\bigr) =
\frac16$，因为 $C \cap \intoo{\frac13}{\frac23} = \emptyset$ 且
$\frac13, \frac23 \in C$。因此 $\max_{y\in\intcc{0}{1}} d(y, C)
= \frac16$，且恰在 $\frac12$ 处取得。

\textbf{24.} 端点构成 $C$ 的一个可数稠密子集
（第 12 题）：将其排成单一序列 $(e_j)_{j\geq1}$，
即 $C$ 中的一个序列。其子列极限全都落在 $C$ 中（$C$
为闭集）。反之，固定 $x \in C$：对每个 $n$，
$C_m$ 中包含 $x$ 的线段（$m \geq n$），其端点
位于半径为 $3^{-m} \leq 3^{-n}$、中心为 $x$ 的邻域内，故无穷多个不同端点
落在半径为 $3^{-n}$、中心为 $x$ 的邻域内；取指标 $j_1 < j_2
< \dots$，使 $\abs{e_{j_n} - x} \leq 3^{-n}$，得一子列
收敛于 $x$。故 $(e_j)$ 的子列极限全体
恰为 $C$ --- 一个序列聚集于不可数多个点，
这正是收敛序列的另一极端，后者
的聚点集为单点集。

\textbf{25.} (i) 构造用到：闭集在任意交下的稳定性（保证
$C$ 为闭集的存在性），
紧性定理（\cref{thm:b1:topology:compact}）
（第 2、22、23 题），以及闭性与闭包的序列刻画
（紧集套论证，见
\cref{ex:b1:topology:nested}，以及第 23 题）。(ii) 四组
配对：长度为零却不可数（第 3、9 题）；闭集
却内部为空（第 4 题）；完备集 --- 无孤立
点 --- 却完全不连通（第 11、13 题）；
可忽略却具有 $C + C = \intcc{0}{2}$（第 16 题）。(iii)
第 20 题中的满射 $h$，经改造为连续且不减的
函数，即成为连续函数理论中的魔鬼阶梯；而在
本科三年级分册的测度论中，
$C$ 是说明“可忽略”并不意味着“可数”的标准例子，
而其较“胖”的近亲（第 22 题）则将“无处稠密”与
“可忽略”区分开来。
\end{solution}
